\subsection{Pointed-associator rigidity for \texorpdfstring{$\sha{R8_3_65}$}{R8-3-65}}
\label{sec:r65-pointed-associator-next}
This refinement removes categorical branches of the connected-algebra
search, while leaving the fusion ring unresolved.

\begin{proposition}
In any complex categorification of $\sha{R8_3_65}$, the pointed subcategory
has nonzero class $k\in H^3(C_4,\C^\times)\cong\Z/4$.
Neither of the two rank-two based module actions can be categorified.
Either rank-four action requires $k=2$; if $k=2$, the subgroup algebra
$\one\oplus g^2$ has a Morita dual with the same Grothendieck ring
$\sha{R8_3_65}$ and pointed class $2$.
\end{proposition}
\begin{proof}
If $k=0$, the actual subgroup algebra for $H=\langle g^2\rangle$ exists.
Its four simple right modules arise from the four free right $H$-orbits.
Each has two left $H$-linearizations, since the possible projective
multiplier is trivializable by $H^2(C_2,\C^\times)=0$.
The bimodule dual has four invertibles forming
$\widehat H\times(C_4/H)\cong V_4$, and a free orbit of four simples of
dimension $d=(11+5\sqrt5)/2$, where $d^2=1+11d$.
This argument uses a subgroup algebra, without assuming a central or
Tannakian lift.

For a free noninvertible orbit over an abelian group $G$ of order four,
write $\rho g=\phi(g)\rho$ and
$\rho^2=u+(\sum_h n_hh)\rho$, with $\sum_hn_h=11$.
Associativity forces $\phi^2=1$ and invariance of $n_h$ under
$\{g\phi(g)^{-1}\}$. If $\phi\ne1$, this subgroup has order two and would
make the coefficient sum even. Thus $\phi=1$.
Duality and the odd coefficient sum permit a translate with $\rho^*=\rho$,
$\rho^2=1+N\rho$ and $N=N^*$.
For $V_4$ the complete possibilities are the $27$ unordered partitions of
$11$ into four parts. For a character parameter $m=\chi(N)$,
\[
 \operatorname{Norm}(f)=16(m^2+4),\qquad
 \operatorname{Norm}(\FPdim(\mathcal D))=2000.
\]
The two-FP-value Galois/pseudo-unitary reduction applies to this dual.
Integrality of its distinguished central dimension $D/f$, taking norms in
the quadratic compositum, implies $125/(m^2+4)\in\Z$.
Hence $|m|=1$ or $11$. Only the partitions
$(0,0,0,11)$, $(0,0,5,6)$ and $(2,3,3,3)$ remain. They are excluded,
respectively, by the general rank-two classification, the repository's
unrestricted rank-four dimension bound, and the exact identification with
the already excluded $\sha{R8_3_1453}$.
The accepted rank-four and $\sha{R8_3_1453}$ results are explicit inputs to
this deduction. Thus $k\ne0$.

In a rank-two action each invertible fixes every module simple, so its
scalar module associators would trivialize $k$. In a rank-four action the
stabilizers are $C_2$, so restriction of $k$ to $H^3(C_2,\C^\times)$ is zero.
The restriction map is reduction modulo two, forcing $k=2$.
Conversely, for $k=2$ the subgroup algebra exists. Its dual has an
order-four invertible group and the same free noninvertible orbit.
The complete $V_4$ comparison above and the corresponding $21$ selfdual
$C_4$ coefficient vectors leave only the ring $\sha{R8_3_65}$.
The rank-four Morita module on the dual side forces its pointed class to
be $2$ as well.
\end{proof}
The categorical cuts reduce the four quadratic rings' aggregate necessary
lists to $144$ algebra objects, $178$ induction factors and $30$ full-action
classes. For $\sha{R8_3_65}$ the remaining action ranks are $4,4,8,8$.
The conditional subgroup dual supplies neither an associator for this ring
nor an equivalence between the dual and the original category.
The complete finite coefficient domains, theorem dependencies and source
hashes are recorded in
\path{evidence/rank8_next_20261021/quadratic/POINTED_COCYCLE.md} and
\path{pointed_dual.py} in that directory, with an independent categorical
review in
\path{evidence/rank8_next_20261021/semion_fermion/construction/R65_ADDENDUM_REVIEW.md}.

\subsection{Three remaining center patterns for \texorpdfstring{$\sha{R8_3_1154}$}{R8-3-1154}}
\label{sec:r1154-three-center-patterns}

The following arguments concern necessary induction and fusion data in the
pseudo-unitary Drinfeld center.  They do not construct a center, a fusion
category, or a categorification of the source ring.  In particular,
$\sha{R8_3_1154}$ remains unresolved.  The previously established reductions to
this pseudo-unitary center problem are used with their original scope.

Write $A_0=\mathbf1,A_1,\ldots,A_4$ for the five distinguished central
simples, in the fixed character order of the retained induction frame, and
write $M$ for the forgetful multiplicity matrix.  All five distinguished
objects are self-dual.  Consequently their fusion matrices are real
symmetric, and scalar products of their fusion-word vectors are the exact
vacuum Verlinde moments.  Number the ten pair words by
$11,12,13,14,22,23,24,33,34,44$.

\begin{proposition}[Joint fourth-moment fibers]
Of the $292$ restriction patterns left by the earlier coordinate bounds,
exactly $224$ admit the required joint pair-product fiber.  Up to permutation
of identical restriction columns, these patterns have exactly $2\,904$
nonnegative integral pair-product completions.
\end{proposition}
\begin{proof}
The previously computed minimum-norm solution $X_0$ to the forgetful
pair-product equations has rank-one residual Gram matrix
\[
 G-X_0^{\mathsf T}X_0=\alpha h^{\mathsf T}h,
 \qquad h=(1,0,1,-1,0,0,0,1,-1,1).
\]
Every completion therefore has the form $X=X_0+t h$, with
$Mt=0$ and $t^{\mathsf T}t=\alpha$.  The allowed integral first coordinate
of each row gives a finite interval.  The common rational denominator is
$205$.  For each restriction type, enumerate all histograms in this
interval with the prescribed number of copies; then impose the linear
restriction equations and the exact squared norm jointly.  Nonnegative
excess costs relative to the nearest allowed coordinate give a finite,
exhaustive transfer computation.  Every retained histogram is checked
against all seven non-unit restriction coordinates and all ten
pair-product integrality and positivity conditions.  An independent
all-seven-coordinate transfer count reproduces the $68$ empty fibers,
$224$ nonempty fibers, and $2\,904$ histograms.
\end{proof}

\begin{proposition}[A forced distinguished action]
The action of $A_2$ excludes $216$ of these $224$ patterns, and $2\,830$
of the $2\,904$ pair-product completions.  The eight remaining patterns
are
\[
 3993,3994,6970,6971,8765,8766,8771,8772,
\]
with $74$ pair-product completions in total.
\end{proposition}
\begin{proof}
Adjoin the eight rows of $M$ to the vacuum, the four distinguished
simple vectors, and the ten pair-word vectors.  Their span $V$ has rank
nine.  It is invariant under $N_{A_2}$.  On pair words this follows either
from the exact fifth/sixth-moment Schur complement or from the vanishing
of the triple residual coefficient on every word containing $2$; on the
forgetful rows it is the induction projection formula.  This invariance
was independently checked using the full $23$-generator Schur identity.

The characteristic polynomial on $V$ is
\[
 x^4(x+1)(x-3)(x+5)(x^2-270x+225).
\]
The complete modular eigenvalue multiset shows that the complementary
action is either zero or minus an orthogonal projector of rank three.
In the first case the full matrix is forced, so any nonintegral entry
excludes the completion.  Equivalently, its remaining trace of the square
is zero, which forces the symmetric complementary action to vanish.

In the second case write $N_{A_2}=F-E$, where $F$ is its fixed action
on $V$.  There are four copies of high restriction type $10$, and
$F_{ii}=67/4$ on each.  Integrality and $0\leq E_{ii}\leq1$ force
$E_{ii}=3/4$ at these four entries.  They exhaust $\operatorname{Tr}(E)=3$,
so every other row and column of $E$ vanishes.  Orthogonality to the
nonzero constant restriction row on this block gives
\[
 E=I_4-\tfrac14 J_4
\]
there.  Thus this case also has a uniquely forced full $A_2$ matrix.
Checking its entries gives the stated counts.  This reasoning and the
complete $2\,904$-histogram domain were independently reproduced.
\end{proof}

\begin{proposition}[Single-support induction]
The four patterns $3993,6970,8765,8771$ are impossible.  The remaining
four patterns have $37$ pair-product completions.
\end{proposition}
\begin{proof}
Let $g=X_3$ be the nontrivial invertible source simple.  All five low
restriction columns have zero $g$ coordinate.  In each of the four named
patterns, the only occupied high type with a nonzero $g$ coordinate is
one copy of type $15$, whose $g$ coordinate is two.  Hence $I(g)=2Z$
for a single central simple $Z$.  But
\[
 \operatorname{Tr}(\theta_{I(g)})
 =\operatorname{Dim}(\mathcal C)\nu_1(g)=0,
\]
whereas the left side equals $2\theta_Z\operatorname{FPdim}(Z)\ne0$.
This uses no twist-order guess or numerical feasibility test.
\end{proof}

\begin{proposition}[The last pattern in the first moment group]
Pattern $3994$ is impossible.  Exactly three necessary restriction
patterns remain: $6971,8766,8772$, with $36$ pair-product completions.
\end{proposition}
\begin{proof}
Pattern $3994$ has one pair-product completion.  Adjoining the forgetful
rows to the known words leaves a rank-one triple residual.  The four
copies of type $10$ have the same complete distinguished joint spectral
tuple $(1,-1,-1,0)$, so every vacuum fusion word is constant on their
forced difference subspace.  Exhausting the resulting finite integral
triple fiber gives exactly one histogram.

The vacuum-word span now has rank ten.  Exact moments through degree eight
verify that it is invariant under all four distinguished actions; it is
the full cyclic span, because there are exactly ten distinct joint
spectral tuples and every vacuum modular weight is positive.  On the
unique high-type-$0$ simple, the fixed cyclic part of $N_{A_4}$ has
\emph{diagonal entry $-2$}.  Its complementary spectrum consists of six
zeros and one eigenvalue $-2$.  The complementary operator is therefore
negative semidefinite, so the actual diagonal entry is at most $-2$.
This contradicts nonnegativity of a fusion matrix.  The unique triple
fiber, degree-eight identities, complete complementary spectrum, and
negative diagonal were independently checked.
\end{proof}

\begin{center}
\begin{tabular}{lrr}
Necessary center constraint & Patterns left & Pair completions left\\
\hline
Earlier coordinate bounds & $292$ & ---\\
Joint pair-product fiber & $224$ & $2\,904$\\
Forced $A_2$ action & $8$ & $74$\\
Single-support induction trace & $4$ & $37$\\
Full cyclic negative diagonal & $3$ & $36$
\end{tabular}
\end{center}

The three survivors are necessary center-data patterns, not three
categorifications.  Further finite triple-fiber and bounded-action tests
are recorded separately as ongoing necessary-condition work.  They do
not alter the unresolved status of the source ring.

\paragraph{Separate exact evidence.}
The accepted construction and scoped refinements, with normal and
optimized-Python receipts and independent reviews, are recorded in
\begin{quote}\small\path{evidence/rank8_next_20261021/}.\end{quote}
Their publication overlay preserves
the earlier ledgers and source-bound proof files.
